\section{The $2\frac{2}{3}$-Approximation Algorithm}

The algorithm proposed by authors improves the generic shortest superstring
algorithm by refining the selection of cycle representatives before the second
cycle cover is computed. The key idea is to choose representative strings so
that the overlap between any two representatives is sufficiently small. As a
result, the total overlap lost $OV$ during the cycle-breaking phase is reduced,
leading to a better approximation ratio.

\subsection{Choosing Cycle Representatives}
We must select a representative string $t_C$ for each cycle that satisfies the
superstring containment conditions and whose infinite periodic extension is
aligned with the corresponding optimal rotation from Lemma
\ref{lemma:overlap_rotation}.

\begin{lemma} \label{lemma:representative_choice}
For any cycle $C = s_{i_1}, \dots, s_{i_r}, s_{i_1} \in \mathcal{C}_S$, there
exists a string $t_C$ such that for some index $j$:
\begin{enumerate}
    \item $t_C$ contains the string $\langle s_{i_{j+1}}, \dots, s_{i_r},
    s_{i_1}, \dots, s_{i_j} \rangle$.
    \item $t_C$ is contained in the string $\langle s_{i_j}, \dots, s_{i_r},
    s_{i_1}, \dots, s_{i_{j-1}}, s_{i_j} \rangle$.
    \item $(t_C)_\infty = \overrightarrow{\langle s_{i_1}, \dots, s_{i_r}
    \rangle}_\infty$.
\end{enumerate}
\end{lemma}

\begin{proof}
Order the equivalent cyclic strings $\langle s_{i_1}, \dots, s_{i_r} \rangle,
\dots, \langle s_{i_r}, s_{i_1}, \dots, s_{i_{r-1}} \rangle$ according to their
first appearances in the infinite optimally rotated string
$\overrightarrow{\langle s_{i_1}, \dots, s_{i_r} \rangle}_\infty$. This ordering
is unique.

Suppose that the string $\langle s_{i_{j+1}}, \dots, s_{i_r}, s_{i_1}, \dots,
s_{i_j} \rangle$ appears first in this ordering. Let $t$ be the prefix of
$\overrightarrow{\langle s_{i_1}, \dots, s_{i_r} \rangle}_\infty$ that ends
exactly at $\langle s_{i_{j+1}}, \dots, s_{i_r}, s_{i_1}, \dots, s_{i_j}
\rangle$ (inclusive). Consequently, $t$ is naturally contained in $\langle
s_{i_j}, \dots, s_{i_r}, s_{i_1}, \dots, s_{i_{j-1}}, s_{i_j} \rangle$. If it
were not, the string $\langle s_{i_j}, \dots, s_{i_r}, s_{i_1}, \dots,
s_{i_{j-1}} \rangle$ would have to appear before $\langle s_{i_{j+1}}, \dots,
s_{i_r}, s_{i_1}, \dots, s_{i_j} \rangle$ in the infinite string, which
fundamentally contradicts our initial choice of the earliest appearing string.

The chosen string $t$ is designated as our representative $t_C$. This selection
process can be executed in polynomial (and in fact, linear) time. From Lemma
\ref{lemma:overlap_rotation} we can get $\overrightarrow{\langle s_{i_1}, \dots,
s_{i_r} \rangle}_\infty$ in linear time to $|\langle s_{i_1}, \dots, s_{i_r}
\rangle|$. Then we only need to find which string appears first which can be
easily done by looking which concatenation point appears first in rotated
string.
$$
\langle s_{i_1}, \dots, s_{i_r} \rangle = \textit{pref}(s_{i_1}, s_{i_2})
\textit{pref}(s_{i_2}, s_{i_3}) \dots \textit{pref}(s_{i_{r-1}}, s_{i_r})
s_{i_r}
$$
\end{proof}

\subsection{Algorithm Outline}
With the representative strings correctly defined we can modify steps 3 and 6 of
generic algorithm and obtaine improved version:

\begin{enumerate}
    \item \textbf{Initial Distance Graph:} Construct the distance graph $G_S$
    for the input set $S$.
    \item \textbf{First Cycle Cover:} Find a minimum weight cycle cover
    $\mathcal{C}_S$ on the graph $G_S$.
    \item \textbf{Optimal Representatives:} For each cycle $C \in
    \mathcal{C}_S$, compute and choose the representative string $t_C$ exactly
    as described in Lemma \ref{lemma:representative_choice}.
    \item \textbf{Second Distance Graph:} Let $T$ be the set of all chosen
    strings $t_C$. Construct the distance graph $G_T$.
    \item \textbf{Second Cycle Cover:} Find a minimum weight cycle cover
    $\mathcal{C}_T$ on the graph $G_T$.
    \item \textbf{Strategic Cycle Breaking:} For each cycle of $\mathcal{C}_T$,
    break the cycle by explicitly deleting an edge that goes from a string to a
    string of equal or larger period. This creates a superstring of the elements
    within the cycle.
    \item \textbf{Final Concatenation:} Arbitrarily concatenate the resulting
    strings to produce the final superstring $\tilde{s}$ for the set $S$.
\end{enumerate}

\subsection{Approximation Ratio Analysis}
It is crucial to note that small cycles in $\mathcal{C}_T$ are not given any
special structural treatment. Instead, the algorithm relies entirely on the
precise condition in Step 6: cutting an edge that leads to a string with an
equal or larger period.

\begin{theorem} \label{thm:2_67_ratio}
The length of the final superstring $\tilde{s}$ produced by the algorithm
satisfies $|\tilde{s}| \le 2\frac{2}{3} opt(S) \approx 2.67 opt(S)$.
\end{theorem}

\begin{proof}Let $OV$ denote the total overlap lost by cutting the edges during Step 6. Because the representative strings comprising the set $T$ are mutually inequivalent, we can apply the bounds provided by the Overlap-Rotation Lemma alongside the established period-weight corollaries. Moreover, since Step 6 always cuts an edge leading to a string of equal or larger period, every such edge satisfies $\textit{period}(t_{C}) \le \textit{period}(t_{C'})$ for the pair of representatives it connects, which is exactly the condition required to invoke the tighter $\frac{2}{3}$ bound of the Overlap-Rotation Lemma rather than the general $period(s) + period(\alpha)$ bound.
\begin{equation}
OV \le \frac{2}{3} \sum_{C \in \mathcal{C}_S} \textit{period}(t_C) = \frac{2}{3} \sum_{C \in \mathcal{C}_S} w(C) = \frac{2}{3} CYC(G_S) \le \frac{2}{3} opt(S)
\end{equation}

From Lemma \ref{lemma:optT} we know that the weight of the cycle cover
$CYC(G_T)$ cannot exceed $2opt(S)$, we can sum these components to find the
strict upper bound on the length of the final superstring $\tilde{s}$:
\begin{equation}
|\tilde{s}| = CYC(G_T) + OV \le 2opt(S) + \frac{2}{3}opt(S) = 2\frac{2}{3}opt(S)
\end{equation}
This concludes the proof, guaranteeing the $2\frac{2}{3}$ approximation ratio.
\end{proof}

\subsection{Computational Complexity}
Let $n = |S|$ and let $L = \sum_{s \in S} |s|$ denote the total input length.
Steps 1 and 4 each build a distance graph on at most $n$ vertices. Steps 2 and 5
each compute a minimum weight cycle cover using the Hungarian algorithm, which
runs in $O(n^3)$ time \cite{kuhn1955hungarian, edmondskarp1972, tomizawa1971techniques}. By Lemma
\ref{lemma:representative_choice}, computing the representative $t_C$ for a
cycle $C$ takes time linear in $|\langle s_{i_1}, \dots, s_{i_r} \rangle|$.
Since the cycles of $\mathcal{C}_S$ partition $S$, Step 3 takes time $O(L)$
overall. Steps 6 and 7 are also linear in the size of the strings being
concatenated, so they take $O(L)$ time as well. Since the two cycle cover
computations take $O(n^3)$ time and all remaining steps take $O(L)$ time, the
algorithm runs in $O(n^3 + L)$ time.
